\documentclass[margin=2mm]{standalone}

\usepackage[utf8]{vietnam}
\usepackage{pgfplots}
\pgfplotsset{compat=1.18}
\usepackage{xcolor}
\usepackage{tikz}

\definecolor{CTUblue}{rgb}{0.121568627450,0.3607843137,0.662745098}
\definecolor{CTUcyan}{rgb}{0,0.6862745098,0.937254902}

% ===== Gaussian pdf (mu, sigma) =====
\pgfmathdeclarefunction{gauss}{2}{%
	\pgfmathparse{1/(#2*sqrt(2*pi))*exp(-((x-#1)^2)/(2*#2^2))}
}

\begin{document}
	\begin{tikzpicture}
		\begin{axis}[
				width=8cm, height=7cm,
			domain=-6:6, samples=100,
			axis lines=box,
			CTUblue,
			ticklabel style={font=\small}
			]
			
			% ---- Ba Gaussian ----
			\def\mA{-2} \def\sA{1}
			\def\mB{ 0} \def\sB{.6}
			\def\mC{ 2} \def\sC{1}
			
			% ---- Trọng số ----
			\pgfmathsetmacro{\wA}{1/3}
			\pgfmathsetmacro{\wB}{1/3}
			\pgfmathsetmacro{\wC}{1/3}
			
			% ---- Tính barycenter theo Wasserstein-2 ----
			% Với 1D Gaussian, W2 barycenter có:
			% μ* = ∑ w_i μ_i
			% σ* = ∑ w_i σ_i
			\pgfmathsetmacro{\muW}{\wA*\mA + \wB*\mB + \wC*\mC}
			\pgfmathsetmacro{\sigW}{\wA*\sA + \wB*\sB + \wC*\sC}
			
			% ---- Các đường PDF riêng ----
			\addplot[CTUblue, dashed, line width=1pt] {gauss(\mA,\sA)};
			\addplot[CTUblue, dashed, line width=1pt] {gauss(\mB,\sB)};
			\addplot[CTUblue, dashed, line width=1pt] {gauss(\mC,\sC)};
			
			% ---- Barycenter theo Wasserstein-2 ----
			\addplot[line width= 2pt, fill=CTUcyan, opacity=.4,CTUcyan]
			{gauss(\muW,\sigW)};
			
		\end{axis}
	\end{tikzpicture}
\end{document}
